\cleardoublepage
\chapter{HASIL UJI STATISTIK DAN PROTOKOL PERCOBAAN}
\label{lamp:statistik}

Lampiran ini memuat dua hal yang dirujuk Bab~\ref{chap:evaluasi} tetapi terlalu rinci untuk badan laporan: hasil uji statistik per-\emph{fold} beserta ukuran efeknya, serta daftar protokol percobaan yang ditulis sebelum hasilnya dilihat.

\section{Hasil Uji Statistik Lengkap}
\label{lamp:uji-statistik}

Seluruh uji dijalankan atas enam \emph{fold} validasi silang lintas-pasien, dengan \emph{fold} sebagai satuan analisis, mengikuti prosedur pada Subbab~\ref{subsec:bab6-signifikansi-retrieval}. Nilai yang dibandingkan adalah MRR pada lengan penelusuran produksi. Ukuran efek yang dilaporkan adalah Cohen's $d_z$ untuk rancangan berpasangan.

% Sumber angka: results/retrieval_realcases_kb12/crossfold.json, lengan bm25.
% Dipilih karena konfigurasi_efektif-nya cocok penuh dengan config.yaml produksi:
% chunk 900, overlap 120, all-MiniLM-L6-v2, top_k 5, collection diabetes_kb.
% Dihitung ulang dari per_fold, bukan disalin dari ringkasan. Wilcoxon dua sisi, n=6 fold.
\begingroup
\setlength{\LTcapwidth}{\textwidth}
\small
\captionsetup{font=normalsize}
\begin{longtable}{@{}lrrrrrr@{}}
\caption{Hasil uji berpasangan lintas enam \emph{fold} pada lengan penelusuran produksi. Nilai MRR, uji Wilcoxon dua sisi, dan ukuran efek Cohen's $d_z$.}
\label{tab:uji-statistik} \\
\toprule
\textbf{Perbandingan} & \textbf{Perlakuan} & \textbf{Kendali} & \textbf{Selisih} & \textbf{SD} & \textbf{$d_z$} & \textbf{$p$} \\
\midrule
\endfirsthead

\multicolumn{7}{@{}l}{\small\emph{Tabel \thetable{} (lanjutan)}} \\
\toprule
\textbf{Perbandingan} & \textbf{Perlakuan} & \textbf{Kendali} & \textbf{Selisih} & \textbf{SD} & \textbf{$d_z$} & \textbf{$p$} \\
\midrule
\endhead

\bottomrule
\endlastfoot

\multicolumn{7}{@{}l}{\textit{(a) Himpunan divergen: kondisi terkini $\neq$ kondisi masa depan}} \\

\midrule

Terkondisi (regresi) & 0,284 & 0,190 & $+0{,}095$ & 0,024 & 3,99 & 0,031 \\

Terkondisi (pengklasifikasi) & 0,312 & 0,190 & $+0{,}122$ & 0,017 & 7,33 & 0,031 \\

\midrule

\multicolumn{7}{@{}l}{\textit{(b) Himpunan natural: distribusi jendela apa adanya}} \\

\midrule

Terkondisi (regresi) & 0,714 & 0,700 & $+0{,}014$ & 0,026 & 0,54 & 0,312 \\

Terkondisi (pengklasifikasi) & 0,663 & 0,700 & $-0{,}036$ & 0,040 & $-0{,}91$ & 0,156 \\
\end{longtable}
\endgroup

\vspace{2pt}
{\tabnotestyle Kendali adalah RAG standar yang membentuk kueri dari kondisi terkini. Selisih dan SD dihitung atas selisih per-\emph{fold}. Kedua baris pada (a) unggul pada seluruh enam \emph{fold}; $p=0{,}031$ merupakan nilai terkecil yang dapat dicapai enam pasangan. Pada (b) tidak satu pun perbandingan mencapai signifikansi, sementara jalur pengklasifikasi bahkan berarah negatif dan tidak unggul pada satu \emph{fold} pun. Dibandingkan terhadap jalur regresi, bukan terhadap RAG standar, kekalahan pengklasifikasi pada himpunan natural terjadi pada \textbf{seluruh enam \emph{fold}} ($-0{,}050$; $d_z=-1{,}52$; $p=0{,}031$), sehingga ia satu-satunya hasil signifikan pada bagian (b).\par}


Tiga hal perlu dibaca bersama tabel di atas.

\textbf{Nilai $p=0{,}031$ adalah batas bawah rancangan ini}, bukan hasil yang nyaris tidak lolos. Pada uji Wilcoxon dua sisi dengan enam pasangan, keadaan paling ekstrem yang mungkin terjadi, yakni keenam selisih searah, menghasilkan tepat $2/2^{6}=0{,}031$. Kedua baris pada himpunan divergen sudah berada pada batas itu.

\textbf{Nilai $d_z$ yang besar mencerminkan keajekan, bukan besar manfaat klinis.} Penyebutnya adalah simpangan baku selisih antar-\emph{fold}, sehingga efek yang sederhana besarannya tetap menghasilkan $d_z$ besar apabila ia bekerja pada tiap \emph{fold} tanpa kecuali. Besar manfaatnya dibaca dari kolom selisih mutlak.

\textbf{Himpunan natural tidak menunjukkan keunggulan yang dapat diandalkan}. Hal itu dilaporkan apa adanya. Pada jalur pengklasifikasi, arahnya bahkan negatif. Ini konsisten dengan penjelasan pada Bab~\ref{chap:evaluasi}: pada mayoritas jendela, kondisi terkini dan kondisi masa depan sama, sehingga tidak ada yang dapat diperbaiki oleh pengondisian.

\section{Daftar Protokol Percobaan}
\label{lamp:protokol}

Tiap protokol ditulis sebelum hasil percobaannya dilihat, memuat dugaan, ambang keputusan, dan aturan penafsiran. Kolom tanggal diambil dari riwayat berkas kerja, yakni saat protokol pertama kali dicatatkan, sehingga urutannya terhadap pelaksanaan percobaan dapat diperiksa ulang. Ringkasan dugaan dan ambangnya disajikan langsung pada tabel ini agar lampiran ini dapat dibaca berdiri sendiri, tanpa menuntut akses ke berkas kerja penelitian.

% Tabel MANDIRI: memuat ringkasan dugaan dan ambang, bukan lintasan berkas kerja.
% Berkas protokolnya sendiri adalah dokumen kerja internal dan tidak dipublikasikan.
% Tanggal diambil dari riwayat berkas (commit pertama tiap protokol).
\begin{longtable}{@{}p{4.6cm}p{1.85cm}p{6.6cm}@{}}
\caption{Protokol percobaan yang ditulis sebelum hasilnya dilihat, beserta ambang keputusan yang ditetapkan di muka.}
\label{tab:protokol} \\
\toprule
\textbf{Percobaan} & \textbf{Tanggal} & \textbf{Dugaan dan ambang yang ditetapkan di muka} \\
\midrule
\endfirsthead
\multicolumn{3}{@{}l}{\footnotesize\textit{Tabel~\ref{tab:protokol} (lanjutan)}} \\
\toprule
\textbf{Percobaan} & \textbf{Tanggal} & \textbf{Dugaan dan ambang yang ditetapkan di muka} \\
\midrule
\endhead
\midrule
\multicolumn{3}{r@{}}{\footnotesize\textit{bersambung}} \\
\endfoot
\bottomrule
\endlastfoot

Sensitivitas konstanta peluruhan IOB dan COB & 11-08-2026 & Mutu prediksi tidak berubah nyata terhadap variasi $\tau$; bila berubah, konstanta harus disetel dan bukan diambil dari pustaka \\[3pt]
Varian kueri berurutan kondisi & 11-08-2026 & Urutan penyebutan kondisi pada kueri tidak mengubah peringkat dokumen \\[3pt]
Jangkauan varian kueri dan bentuk kueri aplikasi & 11-08-2026 & Bentuk kueri yang dipakai aplikasi setara dengan bentuk yang dipakai evaluasi \\[3pt]
\emph{Gradient boosting} sebagai pembanding ketiga & 11-08-2026 & Selisih yang lebih kecil daripada simpangan baku antar-\emph{fold} tidak boleh dinarasikan sebagai keunggulan \\[3pt]
Perbaikan instrumen evaluasi, putaran kedua & 13-08-2026 & Perbaikan alat ukur dijalankan tanpa melihat pengaruhnya terhadap hasil yang sudah ada \\[3pt]
Kestabilan tiga metrik RAGAS & 13-08-2026 & Metrik yang reratanya bergeser antar-jalan identik tidak dipakai untuk menyimpulkan \\[3pt]
Penghitungan ulang penelusuran dan SMBG dengan prediktor produksi & 16-08-2026 & Mode \emph{standard} dan \emph{oracle} wajib tereproduksi persis; bila tidak, yang berubah adalah jalur pengukurannya \\[3pt]
Strategi pemecahan dokumen sadar-kalimat & 16-08-2026 & Pemisah kalimat diadopsi atas dasar integritas korpus, tanpa menunggu MRR; MRR hanya dilaporkan \\[3pt]
Penyelarasan model kalibrasi konformal dengan model produksi & 16-08-2026 & Angka cakupan akan berubah; arah perubahannya tidak boleh ditetapkan sesudah hasil dilihat \\[3pt]
Uji lengan penelusuran pada validasi silang & 17-08-2026 & Bila penelusuran leksikal sendirian mengungguli penggabungan hibrida, maka leksikal yang diadopsi \\[3pt]
Validasi lengan penelusuran dengan alat ukur bebas kata kunci & 17-08-2026 & Keputusan diambil dari \emph{context recall}; selisih di bawah 0,05 dihitung setara; aturan berlaku meski mencabut komponen yang sudah dipasang \\

\end{longtable}

